\itemhead{第3章 小结与复习 · 第5题（教材p.89）}{XJ-B1-C03-SU-P089-Y1370}
\label{XJ-B1-C03-SU-P089-Y1370}
{\footnotesize 湘教版教材p.89 / PDF第96页 / 第3章 小结与复习 / 第5题}
\field{原题}{函数是描述客观世界变化规律的数学模型。函数的思想方法不仅将贯穿高中数学课程的始终，还会成为你今后学习或用到的数学知识的主旋律。小学到初中所学的数学知识，大都可以放到函数的框架之中；大量的实际问题，可以用函数模型来描述和解决。用函数的观点看数学和其他学科里的许多问题，常有纲举目张的效果。}
\field{参考答案}{示例：匀速运动 $s=vt$、购物 $y=px$、圆面积 $S=\pi r^2$ 均为函数模型。建模时明确自变量、实际定义域、单位、关系式，再用结果解释实际问题。本条为反思性总结，不设唯一数值答案。}
{\small \field{关联知识}{F03 函数求值、解析式与函数建模（教材帮p.80、87、88、89、90、131 / PDF89、96、97、98、99、140）}
\field{方法与易错点}{标清每次输入和输出；按已知函数类型或代换关系建立等价方程。 易错：将 $f(x)$ 当成 $f	imes x$；复合顺序反；换元后沿用错误定义域。}
\field{答案依据}{答案册未提供本栏目的独立对应答案；依据原题、教材概念补写参考回应。开放性总结不设唯一数值答案。}}
